% !TEX program = xelatex
\documentclass[UTF8,zihao=-5]{ctexart}
\usepackage[a4paper,top=6mm,bottom=7mm,left=6mm,right=6mm,columnsep=4.5mm]{geometry}
\usepackage{fontspec}
\usepackage{multicol}
\usepackage{enumitem}
\usepackage{tabularx,array,booktabs}
\usepackage{xcolor}
\usepackage{fancyvrb}
\usepackage{amsmath,amssymb}
\usepackage{titlesec}
\usepackage{microtype}

\IfFontExistsTF{Microsoft YaHei}{\setCJKmainfont{Microsoft YaHei}}{}
\IfFontExistsTF{Arial}{\setmainfont{Arial}}{}
\IfFontExistsTF{Consolas}{\setmonofont{Consolas}}{}

\pagestyle{empty}
\setlength{\parindent}{0pt}
\setlength{\parskip}{0.6pt}
\setlength{\columnsep}{4.5mm}
\setlength{\columnseprule}{0.15pt}
\setlength{\premulticols}{0pt}
\setlength{\postmulticols}{0pt}
\setlength{\multicolsep}{1.5pt}
\linespread{0.86}
\sloppy
\emergencystretch=2em

\titleformat{\section}{\bfseries\large}{ }{0pt}{}[\vspace{-1pt}\titlerule]
\titleformat{\subsection}{\bfseries\normalsize}{ }{0pt}{}
\titlespacing*{\section}{0pt}{3pt}{1.5pt}
\titlespacing*{\subsection}{0pt}{2pt}{0.5pt}
\setlist[itemize]{leftmargin=1.05em,itemsep=0pt,topsep=0.5pt,parsep=0pt,partopsep=0pt}
\setlist[enumerate]{leftmargin=1.25em,itemsep=0pt,topsep=0.5pt,parsep=0pt,partopsep=0pt}
\renewcommand{\arraystretch}{0.9}
\newcolumntype{Y}{>{\raggedright\arraybackslash}X}
\DefineVerbatimEnvironment{Code}{Verbatim}{fontsize=\scriptsize,breaklines=true,breakanywhere=true,baselinestretch=0.82,xleftmargin=0pt}
\newcommand{\ans}[1]{\textbf{Ans: #1}}
\newcommand{\hint}[1]{{\color{gray!35!black}\footnotesize #1}}
\newcommand{\okmark}{\ensuremath{\checkmark}}

\begin{document}
\fontsize{7.8pt}{8.7pt}\selectfont
\begin{center}
{\bfseries\Large EE114 嵌入式系统开卷考试 Cheatsheet}\par
{\scriptsize 覆盖 Lect1--18：8051 → ARM Cortex-M3 / STM32F103 → 外设 GPIO/USART/SPI/I2C/Timer/PWM/ADC/DAC/DMA → CPS / 传感器 / 电路常识}\par
\end{center}
\vspace{-5pt}
\begin{multicols*}{2}
\hint{覆盖 Lect1–18：8051 → ARM Cortex‑M3 / STM32F103 → 外设（GPIO/USART/SPI/I2C/Timer/PWM/ADC/DAC/DMA） → CPS / 传感器 / 电路常识}\par
\vspace{1pt}\hrule\vspace{1pt}
\section*{1. 计算机体系结构基础}
{\scriptsize\setlength{\tabcolsep}{1.8pt}\renewcommand{\arraystretch}{0.88}%
\begin{tabularx}{\linewidth}{@{}*{3}{Y}@{}}
\toprule
\textbf{项} & \textbf{冯·诺依曼 (Von Neumann)} & \textbf{哈佛 (Harvard)} \\
\midrule
程序与数据 & \textbf{存储在同一存储器，统一编址} & 程序ROM、数据RAM \textbf{分开} \\
总线 & 一条 & 指令总线 + 数据总线（可并行） \\
代表 & 8086、早期PC & \textbf{8051、ARM Cortex‑M3} \\
\bottomrule
\end{tabularx}}
\begin{itemize}
\item \textbf{三级流水线顺序}：\textbf{取指 (Fetch) → 译码 (Decode) → 执行 (Execute)}
\item \textbf{CPU基本组成}：ALU、寄存器组、控制器（CU）、PC、IR、内部总线
\item \textbf{总线}：地址总线（单向）、数据总线（双向）、控制总线
\item \textbf{RAM}：易失，存储\textbf{动态加载的程序代码、变量、堆栈}；\textbf{ROM/Flash}：存储启动代码与程序
\end{itemize}
\subsection*{CPS / 嵌入式系统}
\begin{itemize}
\item \textbf{CPS (Cyber‑Physical Systems)}：感知 → 计算 → 通信 → 执行 → 反馈，物理过程与计算紧耦合
\item \textbf{嵌入式系统}：面向应用，软硬件可裁剪，对实时性、功耗、成本敏感
\item \textbf{Transducer}：传感器（物理量→电信号）、执行器（电信号→物理量），需调理电路（放大/滤波/隔离）
\item \textbf{采样定理}：fs ≥ 2 fmax 否则混叠
\end{itemize}
\vspace{1pt}\hrule\vspace{1pt}
\section*{2. 8051 单片机要点}
\subsection*{2.1 资源 \& 体系}
\begin{itemize}
\item 8 位 CISC，\textbf{哈佛结构}；\textbf{机器周期 = 12 个振荡周期}
\item 内部 ROM 4KB（8052: 8KB），内部 RAM 128B（8052: 256B），SFR 区 80H–FFH
\item 数据指针 DPTR=16 bit（DPH/DPL）；堆栈指针 SP 复位值 = 07H
\item PSW 标志位：CY、AC、F0、RS1、RS0、OV、—、P
\end{itemize}
\subsection*{2.2 GPIO P0–P3}
\begin{itemize}
\item \textbf{P0} 内部无上拉，作 I/O 需外接上拉；可复用为地址/数据低 8 位
\item \textbf{P1} 准双向，内部上拉
\item \textbf{P2} 准双向，作高 8 位地址
\item \textbf{P3} 准双向，复用功能：
\end{itemize}
{\scriptsize\setlength{\tabcolsep}{1.8pt}\renewcommand{\arraystretch}{0.88}%
\begin{tabularx}{\linewidth}{@{}*{2}{Y}@{}}
\toprule
\textbf{Pin} & \textbf{功能} \\
\midrule
P3.0 & RXD \\
P3.1 & TXD \\
P3.2 & /INT0 \\
P3.3 & /INT1 \\
P3.4 & T0 \\
P3.5 & T1 \\
P3.6 & /WR \\
P3.7 & /RD \\
\bottomrule
\end{tabularx}}
\subsection*{2.3 中断（5 个中断源）}
{\scriptsize\setlength{\tabcolsep}{1.8pt}\renewcommand{\arraystretch}{0.88}%
\begin{tabularx}{\linewidth}{@{}*{3}{Y}@{}}
\toprule
\textbf{源} & \textbf{向量} & \textbf{优先顺位（默认）} \\
\midrule
/INT0 & 0003H & 1（最高） \\
T0 & 000BH & 2 \\
/INT1 & 0013H & 3 \\
T1 & 001BH & 4 \\
串口 RI/TI & 0023H & 5 \\
\bottomrule
\end{tabularx}}
\begin{itemize}
\item IE = EA / — / ET2 / ES / ET1 / EX1 / ET0 / EX0
\item 中断函数：\texttt{void isr() interrupt n [using m]}
\end{itemize}
\subsection*{2.4 定时器 / 计数器 (T0, T1)}
\begin{itemize}
\item TMOD 高 4 位控 T1，低 4 位控 T0；GATE、C/T、M1、M0
\item 模式：\textbf{0=13bit, 1=16bit, 2=8bit 自动重装, 3=T0 分裂为两个 8bit}
\item TCON：TF1 TR1 TF0 TR0 IE1 IT1 IE0 IT0
\item \textbf{软件定时缺点}：占 CPU、精度差、不能后台运行 ➜ 不可说"没有缺点"
\end{itemize}
\subsection*{2.5 串口}
\begin{itemize}
\item SCON：SM0 SM1 SM2 REN TB8 RB8 TI RI
\item 模式 1（8 位 UART）波特率 = (2\textasciicircum{}SMOD / 32) × T1 溢出率
\item 9600bps @ 11.0592MHz，模式 2 → TH1 = TL1 = 0xFD
\end{itemize}
\vspace{1pt}\hrule\vspace{1pt}
\section*{3. C 语言常考易错点}
\subsection*{3.1 宏定义陷阱（重点！）}
\begin{itemize}
\item 宏是\textbf{纯文本替换}，没有运算优先级 ➜ 必须给参数与整体加括号。
\item \texttt{\#define MIN(x,y) (x)\textless{}(y)?(x):(y)}
\item \texttt{10 * MIN(i,j)} 展开为 \texttt{10*(i)\textless{}(j)?(i):(j)} ≡ \texttt{(10*i\textless{}j)?i:j}
\item i=10,j=15 ⇒ \texttt{100\textless{}15? i:j} ⇒ 取 \textbf{j = 15} ➜ 输出 \textbf{15}
\item \texttt{\#define N 2}, \texttt{\#define Y(n) ((N+1)*n)}, \texttt{z=2*(N+Y(5))}
\item = \texttt{2*(2 + (2+1)*5)} = \texttt{2*(2+15)} = \textbf{34}
\end{itemize}
\subsection*{3.2 其它}
\begin{itemize}
\item \texttt{main()} 在源程序中\textbf{位置任意}（C 不强制）
\item \texttt{\&} 位与 vs \texttt{\&\&} 逻辑与；\texttt{|} vs \texttt{||}；\texttt{==} vs \texttt{=}
\item \texttt{volatile} 修饰中断 / 寄存器 / 外部硬件变量，防止编译器优化
\item C 运算符优先级（高→低）：\texttt{() [] -\textgreater{} .} → 单目 → \texttt{* / \%} → \texttt{+ -} → \texttt{\textless{}\textless{} \textgreater{}\textgreater{}} → 关系 → \texttt{==/!=} → \texttt{\& \textasciicircum{} |} → \texttt{\&\& ||} → \texttt{?:} → 赋值 → 逗号
\item \texttt{sizeof} 是运算符，返回 \texttt{size\_t}
\end{itemize}
\subsection*{3.3 STM32 寄存器三种访问方式}
\begin{Code}
GPIOB->ODR = 0xFFFF;                            // 标准库结构体（最推荐）
*(unsigned int*)(0x40010C0C) = 0xFFFF;          // 直接绝对地址
GPIO_SetBits(GPIOB, GPIO_Pin_All);              // 标准外设库函数
\end{Code}
\hint{\texttt{0xFFFF} 仅 16 位、是 GPIOB 的 ODR；直接指针方式不需要 \texttt{\#define}、但要\textbf{类型转换}才能取写。}\par
\vspace{1pt}\hrule\vspace{1pt}
\section*{4. ARM Cortex‑M3 \& STM32F103}
\subsection*{4.1 内核}
\begin{itemize}
\item 32‑bit ARMv7‑M，\textbf{哈佛结构}，\textbf{三级流水线}（取指‑译码‑执行）
\item Thumb‑2 指令集（16/32 位混合），位带 (bit‑banding)，硬件除法
\item \textbf{NVIC（嵌套向量中断控制器）}集成在内核
\item 寄存器：R0–R12、R13=SP（MSP/PSP）、R14=LR、R15=PC，xPSR
\end{itemize}
\subsection*{4.2 总线 / 存储映射}
\begin{itemize}
\item \textbf{AHB}：高速（≤72MHz）；\textbf{APB2}：高速外设（≤72MHz）；\textbf{APB1}：低速（≤36MHz）
\item 0x0800\_0000：Flash（启动后映射 0x0000\_0000）
\item 0x2000\_0000：SRAM
\item 0x4000\_0000：外设；0xE000\_0000：PPB（NVIC/SysTick）
\item BOOT0/BOOT1 选择启动源：Flash / 系统存储器(Bootloader) / SRAM
\end{itemize}
{\scriptsize\setlength{\tabcolsep}{1.8pt}\renewcommand{\arraystretch}{0.88}%
\begin{tabularx}{\linewidth}{@{}*{4}{Y}@{}}
\toprule
\textbf{BOOT0} & \textbf{BOOT1} & \textbf{启动介质} & \textbf{用途} \\
\midrule
0 & X & 主 Flash & 正常运行用户程序 \\
1 & 0 & 系统存储器 & 内置 Bootloader，可经 USART 等下载 \\
1 & 1 & SRAM & 调试/特殊场景 \\
\bottomrule
\end{tabularx}}
\subsection*{4.3 CMSIS / 库分层}
\begin{itemize}
\item \textbf{CMSIS}：ARM 规定的统一硬件抽象\textbf{标准}（不是强制使用，但厂商遵守）。如 \texttt{SystemInit()}、\texttt{NVIC\_SetPriority()}。
\item \textbf{HAL / 标准外设库}：直接寄存器封装 → 函数接口
\item \textbf{BSP（板级支持包）}：在 HAL 库基础上\textbf{再次封装，隐藏板级硬件细节，给应用层统一接口}（≠ 启动文件，≠ OS 内核）
\item 关键头文件：
\item \texttt{stm32f10x.h} ≈ 8051 的 \texttt{reg52.h}：寄存器、结构体、位定义
\item \texttt{stm32f10x\_conf.h}：外设库\textbf{配置开关}（注释/打开使用的模块）
\end{itemize}
\subsection*{4.4 启动流程}
\begin{enumerate}
\item 上电 → 取栈顶地址（向量表 [0]）
\item 取 Reset\_Handler 入口（向量表 [1]）
\item \texttt{startup\_stm32f10x\_xx.s} → 初始化堆栈 → 调用 \textbf{\texttt{SystemInit()}}（配 PLL/时钟） → 跳转 \texttt{\_\_main} → \texttt{main()}
\end{enumerate}
\vspace{1pt}\hrule\vspace{1pt}
\section*{5. RCC 时钟树}
\begin{itemize}
\item \textbf{HSI}：内部 8 MHz RC，\textbf{复位后默认 SYSCLK 源}
\item \textbf{HSE}：外部晶振，常用 8 MHz；\textbf{LSE}：32.768 kHz（RTC）；\textbf{LSI}：约 40 kHz（IWDG / AWU）
\item \textbf{PLL}：倍频；输入可选 HSI/2 或 HSE/HSE÷2，倍频 ×2\textasciitilde{}×16；最高输出 72 MHz
\item \textbf{SYSCLK 来源}：HSI / HSE / \textbf{PLL}（三选一，\okmark）
\item 经 AHB 预分频 → \textbf{HCLK}（≤72M）
\item → APB1 分频 → \textbf{PCLK1}（≤36M）；默认 PCLK1=36M，TIM2\textasciitilde{}7 时钟 = 72M（APB 分频不为 1 时定时器时钟×2）
\item → APB2 分频 → \textbf{PCLK2}（≤72M）；默认 PCLK2=72M，TIM1/8 时钟 = 72M（APB 分频为 1 时不×2），ADC 由 PCLK2÷2/4/6/8
\item → SysTick：HCLK 或 \textbf{HCLK/8}
\item 任何外设使用前都需先在 RCC\_APBxENR/AHBENR 中\textbf{开时钟}
\end{itemize}
\vspace{1pt}\hrule\vspace{1pt}
\section*{6. STM32 GPIO}
\subsection*{6.1 8 种工作模式（CRL/CRH 中 4 位/引脚：CNF[1:0] + MODE[1:0]）}
\textbf{输入 4 种}：模拟输入 / 浮空输入 / 上拉输入 / 下拉输入\par
\textbf{输出 4 种}：\par
\begin{enumerate}
\item 推挽 (PP) 输出
\item 开漏 (OD) 输出（需外部上拉）
\item 复用推挽 (AF\_PP)
\item 复用开漏 (AF\_OD)
\end{enumerate}
\hint{[注意] \textbf{没有 "浮空输出"}（浮空仅是输入）；只有"推挽/开漏/复用推挽/复用开漏"四种输出。}\par
\subsection*{6.2 速度}
\begin{itemize}
\item 2 MHz / 10 MHz / 50 MHz；高频信号（SPI、USART）选 50 MHz
\end{itemize}
\subsection*{6.3 电流}
\begin{itemize}
\item 单个 I/O 最大灌/拉 \textbf{20 mA}；全部 I/O 总电流 \textbf{≤ 150 mA}
\end{itemize}
\subsection*{6.4 复位状态}
\begin{itemize}
\item 复位后绝大多数 I/O \textbf{= 浮空输入}（JTAG/SWD 引脚例外，PA13/14/15、PB3/4 默认复用）
\item \textbf{IDR 复位值不固定}（随外部电平），\textbf{ODR 复位值 = 0x0000}
\end{itemize}
\subsection*{6.5 关键寄存器}
{\scriptsize\setlength{\tabcolsep}{1.8pt}\renewcommand{\arraystretch}{0.88}%
\begin{tabularx}{\linewidth}{@{}*{3}{Y}@{}}
\toprule
\textbf{寄存器} & \textbf{作用} & \textbf{复位值} \\
\midrule
GPIOx\_CRL / CRH & 配置（每脚 4bit） & 0x44444444 = 全浮空输入 \\
GPIOx\_IDR & 输入数据（只读，bit16‑31 保留） & 不定 \\
GPIOx\_ODR & 输出数据 & 0x00000000 \\
GPIOx\_BSRR & \textbf{写 1} 置位 / 复位（原子操作） & — \\
GPIOx\_BRR & 仅清零 & — \\
GPIOx\_LCKR & 配置锁定 & — \\
\bottomrule
\end{tabularx}}
\vspace{1pt}\hrule\vspace{1pt}
\section*{7. NVIC 中断 \& EXTI \& SysTick}
\subsection*{7.1 NVIC}
\begin{itemize}
\item 256 个向量；STM32F103 有 \textbf{60 个可屏蔽外部中断}
\item \textbf{允许多个中断配置成相同优先级}；但抢占优先级相同时\textbf{不能互相嵌套打断}，响应/子优先级只决定 pending 后的响应顺序
\item 固定异常优先级：\textbf{Reset = -3（最高）}，\textbf{NMI = -2}，\textbf{HardFault = -1}；普通可编程优先级数值越小越高
\item 优先级寄存器 8bit，STM32 只用高 \textbf{4 位}（16 级）→ 5 种分组：
\end{itemize}
{\scriptsize\setlength{\tabcolsep}{1.8pt}\renewcommand{\arraystretch}{0.88}%
\begin{tabularx}{\linewidth}{@{}*{3}{Y}@{}}
\toprule
\textbf{Group} & \textbf{抢占 (Preempt)} & \textbf{子 (Sub)} \\
\midrule
0 & 0 bit & 4 bit \\
1 & 1 bit & 3 bit \\
2 & 2 bit & 2 bit \\
3 & 3 bit & 1 bit \\
4 & 4 bit & 0 bit \\
\bottomrule
\end{tabularx}}
\subsection*{7.2 EXTI（外部中断 / 事件）}
\begin{itemize}
\item 19 条：EXTI0–15（GPIO）、16=PVD、17=RTC Alarm、18=USB Wakeup
\item \textbf{同号 Pin 互斥}：同一时刻 PAx、PBx、PCx 只能选一路接到 EXTIx
\item 共享中断号：EXTI9\_5\_IRQn、EXTI15\_10\_IRQn
\item 触发：上升 / 下降 / 双边沿
\item 配置步骤：开 GPIO+AFIO 时钟 → GPIO 输入模式 → \texttt{GPIO\_EXTILineConfig()} 映射端口到 EXTIx → \texttt{EXTI\_Init()} 设置线/触发 → \texttt{NVIC\_Init()} → 写 \texttt{EXTIx\_IRQHandler()} 并清 pending
\end{itemize}
\subsection*{7.3 SysTick}
\begin{itemize}
\item 24 位\textbf{向下计数}，自动重装
\item CLKSOURCE：1=HCLK（72MHz）/ 0=HCLK/8（9MHz）
\item \texttt{SysTick\_Config(72000)} → 1 ms 中断
\item 常用于 OS Tick 与延时
\end{itemize}
\vspace{1pt}\hrule\vspace{1pt}
\section*{8. USART / UART}
\begin{itemize}
\item 异步：1 起始 + 8/9 数据 + 可选奇偶 + 1/1.5/2 停止
\item 波特率：\texttt{USARTDIV = fCK / (16 × BaudRate)}，\texttt{BRR} 存 USARTDIV 的 mantissa/fraction（例：USART1@72MHz, 115200bps → USARTDIV≈39.0625 → BRR≈0x0271）
\item 关键寄存器：USART\_SR (TXE/RXNE/TC/IDLE/ORE)、DR、CR1 (UE/TE/RE/PCE/PS/M)、BRR
\item \textbf{TTL 0–3.3V} ；\textbf{RS‑232 ±3\textasciitilde{}15V，反相}；\textbf{RS‑485 差分平衡}
\item 引脚：GPIO 复用推挽 50MHz（TX）+ 浮空/上拉输入（RX）
\end{itemize}
\vspace{1pt}\hrule\vspace{1pt}
\section*{9. SPI}
\begin{itemize}
\item 4 线全双工同步：\textbf{SCK / MOSI / MISO / NSS}
\item 4 种模式 by \textbf{CPOL/CPHA}：
\item 0/0 空闲低，第一个上升沿采样
\item 0/1 空闲低，第一个下降沿采样
\item 1/0 空闲高，第一个下降沿采样
\item 1/1 空闲高，第一个上升沿采样
\item STM32F1：\texttt{fSCK = fPCLK / Prescaler}，最小分频为 2。\textbf{SPI1 在 APB2，最高 36 Mbit/s = PCLK2/2}；\textbf{SPI2/3 在 APB1，最高 18 Mbit/s = PCLK1/2}。
\item 寄存器 SPI\_CR1（BR[2:0]分频、CPOL、CPHA、MSTR、SSI、SSM、LSBFIRST、SPE）
\end{itemize}
\vspace{1pt}\hrule\vspace{1pt}
\section*{10. I²C}
\begin{itemize}
\item 两线半双工：SCL + SDA（\textbf{开漏}，必须\textbf{外部上拉电阻}\textasciitilde{}4.7kΩ）
\item 起始：SCL=H 时 SDA ↓；停止：SCL=H 时 SDA ↑
\item 速率：标准 100k / 快速 400k / 高速 3.4M
\item 第 9 bit 为 ACK（接收方拉低 SDA）
\item 7 bit 或 10 bit 地址 + R/W
\item 多主仲裁 + 时钟同步
\end{itemize}
\vspace{1pt}\hrule\vspace{1pt}
\section*{11. 定时器 \& PWM}
\subsection*{11.1 TIM 分类（STM32F103）}
{\scriptsize\setlength{\tabcolsep}{1.8pt}\renewcommand{\arraystretch}{0.88}%
\begin{tabularx}{\linewidth}{@{}*{3}{Y}@{}}
\toprule
\textbf{类型} & \textbf{编号} & \textbf{特性} \\
\midrule
高级 & TIM1, TIM8 & 16 bit, 带死区 / 互补输出 / 刹车 \\
通用 & TIM2–5 & 16 bit, 4 通道，可生成 PWM \\
基本 & TIM6, TIM7 & 仅做计数 / 触发 DAC，无 I/O \\
\bottomrule
\end{tabularx}}
\hint{[\okmark] 通用 \textbf{和} 高级定时器都可输出 PWM；基本定时器不行。}\par
\subsection*{11.2 关键寄存器}
\begin{itemize}
\item TIMx\_PSC（预分频）：实际分频 = PSC+1
\item TIMx\_ARR（自动重装）：周期 = (PSC+1)(ARR+1)/fCK
\item TIMx\_CNT：当前计数
\item TIMx\_CCRx：捕获 / 比较值
\end{itemize}
\subsection*{11.3 PWM}
\begin{itemize}
\item \textbf{频率由 ARR 决定}；\textbf{占空比由 CCR 决定}（占空比 = CCR/(ARR+1)）
\item 模式：
\item \textbf{PWM1（mode 6）}：CNT < CCR 时输出\textbf{有效电平（高）}，≥ 则低
\item \textbf{PWM2（mode 7）}：CNT < CCR 时输出\textbf{低}，≥ 则高
\item 两者相位差 180°；亦可用 OCxREF 极性反相实现，但 PWM1/2 仍有意义（影响极性 + 触发逻辑）
\item 对齐：边沿对齐 / 中央对齐 1/2/3
\end{itemize}
\vspace{1pt}\hrule\vspace{1pt}
\section*{12. ADC / DAC}
\subsection*{12.1 ADC}
\begin{itemize}
\item 12 位 \textbf{逐次逼近 (SAR)}；分辨率 = Vref / 4096（典型 3.3 V → 0.806 mV）
\item 转换流程：\textbf{采样 → 量化 → 编码}
\item 通道：16 个外部 + 内部温度 + Vrefint；规则组（最多 16）+ 注入组（最多 4）
\item 模式：单次 / 连续 / 扫描 / 不连续
\item ADCCLK 最高 14 MHz（由 PCLK2 分频 2/4/6/8 得到）
\item 转换时间 = 采样时间 + 12.5 ADCCLK
\end{itemize}
\subsection*{12.2 DAC}
\begin{itemize}
\item 12 位电压输出，2 通道；可触发产生噪声 / 三角波；可与定时器、DMA 联动
\item 输出缓冲使能时低输出阻抗
\end{itemize}
\subsection*{12.3 DMA}
\begin{itemize}
\item F103：DMA1（7 通道）+ DMA2（5 通道）
\item 方向：M2M / M2P / P2M
\item 循环模式（Circular）：到达末尾自动从头开始
\item 优先级：软件 4 级（低/中/高/极高）+ 硬件（通道号小优先）
\item 配合 ADC/USART/SPI/I²C 减少 CPU 占用
\item 过程：DMA 请求 → 仲裁/响应 → DMA 获得总线并按源/目的地址传输 → 计数到 0 后置标志/中断，CPU 不再逐字节搬运
\end{itemize}
\vspace{1pt}\hrule\vspace{1pt}
\section*{13. 电路 \& 硬件常识}
\begin{itemize}
\item \textbf{三极管开关电路设计参数}：耐压 (Vce(max))、功率 (Pmax)、负载功率、\textbf{β 值}（决定 IB = IL/β · 饱和深度）— 这四项\textbf{都要考虑}
\item \textbf{LDO（线性稳压器）特点}：
\item $\checkmark$ 输出\textbf{纹波小}
\item $\times$ \textbf{效率低}（差压×电流 全部发热）
\item $\times$ \textbf{降压幅度不能大}（差压大则发热大、效率更低）
\item $\times$ 输出\textbf{功率受限}
\item \textbf{共模噪声}：两条线对地同向噪声；\textbf{电池供电可立即排除电源/接地相关的共模噪声}（\okmark 正确）
\item 烙铁头易氧化 → \textbf{使用后给烙铁头上锡保护}（\okmark 正确）
\item \textbf{电路原理图（Schematic）}：用约定符号表示电路结构与工作原理，用于\textbf{设计与分析}（不是 PCB，不"承载"元器件）；PCB 才是支撑体。
\item 共模 vs 差模：滤波电感的方向不同；EMI 设计中区别处理
\end{itemize}
\vspace{1pt}\hrule\vspace{1pt}
\section*{14. 开发环境（MDK / Keil）}
\begin{itemize}
\item 编译链 ARMCC + 链接 + .hex/.bin/.axf
\item 调试器：ST‑Link / J‑Link，接口 SWD（2 线）/ JTAG（4 线）
\item 调试常用操作：\textbf{单步 (Step)}、\textbf{断点 (Breakpoint)}、\textbf{暂停 (Halt) 查看当前 PC}、\textbf{变量加入 Watch 窗}、寄存器/外设窗、内存窗 — \textbf{以上四种都常用} \okmark
\end{itemize}
\vspace{1pt}\hrule\vspace{1pt}
\section*{15. 综合速查表（最常考数值/事实）}
{\scriptsize\setlength{\tabcolsep}{1.8pt}\renewcommand{\arraystretch}{0.88}%
\begin{tabularx}{\linewidth}{@{}*{2}{Y}@{}}
\toprule
\textbf{项目} & \textbf{值} \\
\midrule
STM32F103 最高 SYSCLK & 72 MHz \\
AHB / APB2 最大 & 72 MHz \\
APB1 最大 & 36 MHz \\
GPIO 单引脚最大电流 & 20 mA \\
GPIO 全部最大总电流 & 150 mA \\
GPIO 输出速度选项 & 2 / 10 / 50 MHz \\
SPI 最高速率 & SPI1: \textbf{36 Mbit/s = PCLK2/2}；SPI2/3: \textbf{18 Mbit/s = PCLK1/2} \\
ADC 最高采样 & 1 MSPS (14 MHz ADCCLK) \\
ADC 位数 & 12 bit \\
DAC 位数 & 12 bit \\
三级流水线 & 取指→译码→执行 \\
NVIC 固定最高优先级 & Reset = ‑3，NMI = ‑2，HardFault = ‑1 \\
EXTI 线 & 共 19 条 \\
LSE & 32.768 kHz \\
LSI & ≈ 40 kHz \\
HSI & 8 MHz \\
HSE & 4–16 MHz（典型 8 MHz） \\
SysTick 位数 & 24 bit 向下计数 \\
\bottomrule
\end{tabularx}}
\vspace{1pt}\hrule\vspace{1pt}
\section*{16. 关键寄存器速查}
\subsection*{16.1 8051 SFR}
{\scriptsize\setlength{\tabcolsep}{1.8pt}\renewcommand{\arraystretch}{0.88}%
\begin{tabularx}{\linewidth}{@{}*{3}{Y}@{}}
\toprule
\textbf{寄存器} & \textbf{地址/位} & \textbf{作用/易错点} \\
\midrule
\texttt{P0/P1/P2/P3} & 80H/90H/A0H/B0H & GPIO 端口；P0 无内部上拉；P3 复用 RXD/TXD/INT/T0/T1/WR/RD \\
\texttt{SP} & 81H & 堆栈指针，复位值 07H；压栈先加 1 再写入 \\
\texttt{DPL/DPH} & 82H/83H & 16 位数据指针 \texttt{DPTR}，常用于查表/外部存储器访问 \\
\texttt{PCON} & 87H & 电源控制；\texttt{SMOD} 可使串口波特率加倍 \\
\texttt{TCON} & 88H & \texttt{TF1 TR1 TF0 TR0 IE1 IT1 IE0 IT0}；定时器启停/溢出标志/外部中断触发方式 \\
\texttt{TMOD} & 89H & 高 4 位控 T1，低 4 位控 T0：\texttt{GATE C/T M1 M0} \\
\texttt{TL0/TH0/TL1/TH1} & 8A/8C/8B/8D & 定时器初值/计数值；模式 1 为 16 bit，模式 2 为 8 bit 自动重装 \\
\texttt{SCON} & 98H & \texttt{SM0 SM1 SM2 REN TB8 RB8 TI RI}；串口模式/收发使能/收发中断标志 \\
\texttt{SBUF} & 99H & 串口收发缓冲；写发送，读接收 \\
\texttt{IE} & A8H & \texttt{EA - ET2 ES ET1 EX1 ET0 EX0}；\texttt{EA=1} 后各源使能才有效 \\
\texttt{IP} & B8H & 中断优先级；1=高优先级，0=低优先级 \\
\texttt{PSW} & D0H & \texttt{CY AC F0 RS1 RS0 OV - P}；\texttt{RS1/RS0} 选寄存器组 \\
\texttt{ACC/B} & E0H/F0H & 累加器/乘除辅助寄存器 \\
\bottomrule
\end{tabularx}}
\subsection*{16.2 Cortex-M3 / NVIC / SysTick}
{\scriptsize\setlength{\tabcolsep}{1.8pt}\renewcommand{\arraystretch}{0.88}%
\begin{tabularx}{\linewidth}{@{}*{2}{Y}@{}}
\toprule
\textbf{寄存器} & \textbf{作用/关键位} \\
\midrule
向量表 \texttt{[0] / [1]} & \texttt{[0]} 初始 MSP，\texttt{[1]} Reset\_Handler 入口 \\
\texttt{SCB-\textgreater{}VTOR} & 向量表基地址重定位 \\
\texttt{SCB-\textgreater{}AIRCR} & 优先级分组 \texttt{PRIGROUP}；写入需带 \texttt{VECTKEY=0x5FA} \\
\texttt{NVIC\_ISERx / ICERx} & 使能/失能外部 IRQ \\
\texttt{NVIC\_ISPRx / ICPRx} & 置位/清除 pending \\
\texttt{NVIC\_IPRx} & 中断优先级；STM32F1 只实现高 4 位，数值越小优先级越高 \\
\texttt{SysTick-\textgreater{}LOAD} & 24 位重装值；\texttt{ticks-1} 写入 \\
\texttt{SysTick-\textgreater{}VAL} & 当前计数值；写任意值清零 \\
\texttt{SysTick-\textgreater{}CTRL} & \texttt{ENABLE} 使能，\texttt{TICKINT} 中断，\texttt{CLKSOURCE} 选 HCLK 或 HCLK/8，\texttt{COUNTFLAG} 到 0 标志 \\
\bottomrule
\end{tabularx}}
\subsection*{16.3 RCC / AFIO / EXTI}
{\scriptsize\setlength{\tabcolsep}{1.8pt}\renewcommand{\arraystretch}{0.88}%
\begin{tabularx}{\linewidth}{@{}*{3}{Y}@{}}
\toprule
\textbf{模块} & \textbf{寄存器} & \textbf{作用/关键位} \\
\midrule
RCC & \texttt{RCC\_CR} & HSI/HSE/PLL 开关与就绪：\texttt{HSION/HSIRDY/HSEON/HSERDY/PLLON/PLLRDY} \\
RCC & \texttt{RCC\_CFGR} & 系统时钟与分频：\texttt{SW/SWS} 选 SYSCLK，\texttt{HPRE/PPRE1/PPRE2} 分频，\texttt{PLLSRC/PLLMUL} 配 PLL \\
RCC & \texttt{RCC\_APB2ENR} & 使能 GPIOA-G、AFIO、USART1、SPI1、ADC1/2、TIM1 等 \\
RCC & \texttt{RCC\_APB1ENR} & 使能 TIM2-7、USART2/3、SPI2/3、I2C1/2、DAC、CAN 等 \\
RCC & \texttt{RCC\_AHBENR} & 使能 DMA1/DMA2、SRAM/FLITF、FSMC、SDIO 等 \\
RCC & \texttt{RCC\_BDCR} & LSE/RTC/备份域控制 \\
RCC & \texttt{RCC\_CSR} & LSI 和复位标志 \\
AFIO & \texttt{AFIO\_MAPR} & 引脚重映射；也用于关闭/保留 JTAG/SWD \\
AFIO & \texttt{AFIO\_EXTICR[1..4]} & 选择 EXTIx 连接到 PA/PB/PC...；同号 GPIO 互斥 \\
EXTI & \texttt{IMR/EMR} & 中断屏蔽 / 事件屏蔽 \\
EXTI & \texttt{RTSR/FTSR} & 上升沿 / 下降沿触发 \\
EXTI & \texttt{SWIER} & 软件触发 \\
EXTI & \texttt{PR} & pending 标志；\textbf{写 1 清除} \\
\bottomrule
\end{tabularx}}
\subsection*{16.4 GPIO}
{\scriptsize\setlength{\tabcolsep}{1.8pt}\renewcommand{\arraystretch}{0.88}%
\begin{tabularx}{\linewidth}{@{}*{2}{Y}@{}}
\toprule
\textbf{寄存器} & \textbf{作用/关键点} \\
\midrule
\texttt{GPIOx\_CRL} & 配置 pin0-7；每脚 4 bit = \texttt{CNF[1:0] + MODE[1:0]} \\
\texttt{GPIOx\_CRH} & 配置 pin8-15；复位后通常为 \texttt{0x44444444}（浮空输入） \\
\texttt{GPIOx\_IDR} & 输入数据，只读，读实际引脚电平 \\
\texttt{GPIOx\_ODR} & 输出数据；可读写，但读-改-写不如 BSRR 原子 \\
\texttt{GPIOx\_BSRR} & bit0-15 写 1 置位；bit16-31 写 1 复位；原子操作 \\
\texttt{GPIOx\_BRR} & bit0-15 写 1 复位 \\
\texttt{GPIOx\_LCKR} & 锁定 GPIO 配置，需特定写序列 \\
\bottomrule
\end{tabularx}}
常用模式编码思路：\texttt{MODE=00} 输入；\texttt{MODE=01/10/11} 分别为 10/2/50MHz 输出；输入时 \texttt{CNF=00} 模拟、\texttt{01} 浮空、\texttt{10} 上拉/下拉；输出时 \texttt{CNF=00} 推挽、\texttt{01} 开漏、\texttt{10} 复用推挽、\texttt{11} 复用开漏。\par
\subsection*{16.5 USART / SPI / I2C}
{\scriptsize\setlength{\tabcolsep}{1.8pt}\renewcommand{\arraystretch}{0.88}%
\begin{tabularx}{\linewidth}{@{}*{3}{Y}@{}}
\toprule
\textbf{外设} & \textbf{寄存器} & \textbf{作用/关键位} \\
\midrule
USART & \texttt{USART\_SR} & 状态：\texttt{TXE} 发送数据寄存器空，\texttt{TC} 发送完成，\texttt{RXNE} 接收非空，\texttt{ORE/NE/FE/PE} 错误 \\
USART & \texttt{USART\_DR} & 数据收发寄存器；读清 RXNE，写启动发送 \\
USART & \texttt{USART\_BRR} & 波特率分频值 \\
USART & \texttt{USART\_CR1} & \texttt{UE} USART 使能，\texttt{TE/RE} 收发使能，\texttt{RXNEIE/TXEIE/TCIE} 中断，\texttt{M/PCE/PS} 数据位/校验 \\
USART & \texttt{USART\_CR2} & 停止位 \texttt{STOP}，同步时钟相关 \\
USART & \texttt{USART\_CR3} & \texttt{DMAT/DMAR} DMA 收发，\texttt{RTSE/CTSE} 硬件流控 \\
SPI & \texttt{SPI\_CR1} & \texttt{MSTR} 主机，\texttt{BR[2:0]} 分频，\texttt{CPOL/CPHA} 模式，\texttt{SPE} 使能，\texttt{LSBFIRST}，\texttt{SSM/SSI} 软件片选 \\
SPI & \texttt{SPI\_CR2} & \texttt{TXEIE/RXNEIE/ERRIE} 中断，\texttt{TXDMAEN/RXDMAEN} DMA \\
SPI & \texttt{SPI\_SR} & \texttt{TXE} 可写，\texttt{RXNE} 可读，\texttt{BSY} 忙，\texttt{OVR/MODF} 错误 \\
SPI & \texttt{SPI\_DR} & 收发数据；全双工时写入同时移出/移入 \\
I2C & \texttt{I2C\_CR1} & \texttt{PE} 使能，\texttt{START/STOP/ACK}，\texttt{SWRST} 软件复位 \\
I2C & \texttt{I2C\_CR2} & 外设输入时钟频率 \texttt{FREQ}，中断/DMA 使能 \\
I2C & \texttt{I2C\_OAR1/OAR2} & 自身地址，7/10 位寻址 \\
I2C & \texttt{I2C\_DR} & 数据寄存器 \\
I2C & \texttt{I2C\_SR1/SR2} & 状态：\texttt{SB/ADDR/BTF/RXNE/TXE/AF/BERR/ARLO}；读 SR1 后读 SR2 可清 ADDR \\
I2C & \texttt{I2C\_CCR/TRISE} & SCL 时钟控制 / 最大上升时间 \\
\bottomrule
\end{tabularx}}
\subsection*{16.6 TIM / PWM}
{\scriptsize\setlength{\tabcolsep}{1.8pt}\renewcommand{\arraystretch}{0.88}%
\begin{tabularx}{\linewidth}{@{}*{2}{Y}@{}}
\toprule
\textbf{寄存器} & \textbf{作用/关键位} \\
\midrule
\texttt{TIMx\_CR1} & \texttt{CEN} 使能，\texttt{DIR} 方向，\texttt{CMS} 中央对齐，\texttt{ARPE} ARR 预装载 \\
\texttt{TIMx\_PSC} & 预分频值；实际分频 = \texttt{PSC+1} \\
\texttt{TIMx\_ARR} & 自动重装值；决定周期/频率 \\
\texttt{TIMx\_CNT} & 当前计数值 \\
\texttt{TIMx\_DIER} & \texttt{UIE/CCxIE} 中断使能，\texttt{UDE/CCxDE} DMA 请求 \\
\texttt{TIMx\_SR} & \texttt{UIF/CCxIF} 中断标志；处理后清除 \\
\texttt{TIMx\_EGR} & \texttt{UG} 产生更新事件，让 PSC/ARR 立即装载 \\
\texttt{TIMx\_CCMR1/2} & 通道模式；\texttt{OCxM=110} PWM1，\texttt{111} PWM2；\texttt{OCxPE} CCR 预装载 \\
\texttt{TIMx\_CCER} & \texttt{CCxE} 输出使能，\texttt{CCxP} 输出极性 \\
\texttt{TIMx\_CCRx} & 捕获/比较值；PWM 中决定占空比 \\
\texttt{TIM1/8\_BDTR} & 高级定时器：\texttt{MOE} 主输出使能，死区 \texttt{DTG}，刹车 \texttt{BKE/BKP} \\
\bottomrule
\end{tabularx}}
公式：\texttt{T = (PSC+1)(ARR+1)/TIMxCLK}；边沿对齐 PWM 占空比约 \texttt{CCR/(ARR+1)}。改占空比常用 \texttt{TIM\_SetCompare1/2/3/4()} 或直接写 \texttt{TIMx-\textgreater{}CCRx}。\par
\subsection*{16.7 ADC / DAC / DMA}
{\scriptsize\setlength{\tabcolsep}{1.8pt}\renewcommand{\arraystretch}{0.88}%
\begin{tabularx}{\linewidth}{@{}*{3}{Y}@{}}
\toprule
\textbf{外设} & \textbf{寄存器} & \textbf{作用/关键位} \\
\midrule
ADC & \texttt{ADC\_SR} & \texttt{EOC} 规则转换完成，\texttt{JEOC} 注入转换完成，\texttt{AWD} 模拟看门狗 \\
ADC & \texttt{ADC\_CR1} & 扫描 \texttt{SCAN}，中断 \texttt{EOCIE/JEOCIE}，模拟看门狗配置 \\
ADC & \texttt{ADC\_CR2} & \texttt{ADON} 使能，\texttt{CONT} 连续，\texttt{DMA}，\texttt{ALIGN}，\texttt{EXTSEL/EXTTRIG}，\texttt{SWSTART}，\texttt{CAL/RSTCAL} \\
ADC & \texttt{ADC\_SMPR1/2} & 各通道采样时间 \\
ADC & \texttt{ADC\_SQR1/2/3} & 规则组转换序列与长度 \texttt{L} \\
ADC & \texttt{ADC\_DR} & 规则组转换结果 \\
ADC & \texttt{ADC\_JSQR/JDRx} & 注入组序列/结果 \\
DAC & \texttt{DAC\_CR} & \texttt{ENx} 使能，\texttt{BOFFx} 输出缓冲，\texttt{TENx/TSELx} 触发，\texttt{DMAENx} \\
DAC & \texttt{DAC\_DHRx} & 数据保持寄存器，支持 8/12 位、左/右对齐 \\
DAC & \texttt{DAC\_DORx} & 输出数据寄存器 \\
DMA & \texttt{DMA\_ISR} & 全局/完成/半完成/错误标志：\texttt{GIF/TCIF/HTIF/TEIF} \\
DMA & \texttt{DMA\_IFCR} & 写 1 清 DMA 标志 \\
DMA & \texttt{DMA\_CCRx} & \texttt{EN} 使能，\texttt{TCIE/HTIE/TEIE} 中断，\texttt{DIR} 方向，\texttt{CIRC} 循环，\texttt{PINC/MINC} 地址递增，\texttt{PSIZE/MSIZE} 数据宽度，\texttt{PL} 优先级，\texttt{MEM2MEM} \\
DMA & \texttt{DMA\_CNDTRx} & 传输数量 \\
DMA & \texttt{DMA\_CPARx} & 外设地址 \\
DMA & \texttt{DMA\_CMARx} & 存储器地址 \\
\bottomrule
\end{tabularx}}
配置顺序口诀：先开 RCC → 配 GPIO/复用 → 配外设参数 → 配 NVIC/DMA（若用）→ 清标志 → 使能外设/通道 → 写 ISR 或轮询状态位。\par
\vspace{1pt}\hrule\vspace{1pt}
\section*{17. 例题（完整选项 + 答案）}
\subsection*{17.1 Example.txt 原题整理}
\textbf{1. 单选} 冯·诺依曼结构中，程序指令和数据存储在何处？\par
A. 分别存储在独立的指令存储器和数据存储器中\par
B. 存储在高速缓存 Cache 中，独立访问\par
C. 存储在同一个存储器中，统一编址\par
D. 存储在外部存储设备中，通过总线动态加载\par
\textbf{Ans: C}\par
\textbf{2. 单选} 关于 BSP（板级支持包）的作用，哪项最准确？\par
A. BSP 是硬件抽象层，直接操作寄存器以控制 GPIO、定时器等外设\par
B. BSP 是在 HAL 库基础上的再次封装，隐藏硬件细节并为应用层提供统一接口\par
C. BSP 是 STM32 的启动文件，负责初始化堆栈和调用 main 函数\par
D. BSP 是操作系统内核，负责任务调度和内存管理\par
\textbf{Ans: B}\par
\textbf{3. 单选} \texttt{\#define MIN(x,y) (x)\textless{}(y)?(x):(y)}，\texttt{i=10,j=15; k=10*MIN(i,j);} 的输出是？\par
A. 150.0\par
B. 100.0\par
C. 15.0\par
D. 10.0\par
\textbf{Ans: C}。宏展开为 \texttt{(10*i\textless{}j)?i:j}，得到 15。\par
\textbf{4. 单选} RAM 主要用于：\par
A. 保存部分程序代码\par
B. 动态加载的程序代码\par
C. 保持初始化数值\par
D. 启动程序代码\par
\textbf{Ans: B}\par
\textbf{5. 判断} 烙铁头容易被氧化，所以使用后应该给烙铁头上锡。\par
A. 正确\par
B. 错误\par
\textbf{Ans: A}\par
\textbf{6. 多选} 设计三极管开关电路时，需要考虑的参数有：\par
A. 三极管耐压\par
B. 三极管功率\par
C. 负载功率\par
D. beta 值\par
\textbf{Ans: A B C D}\par
\textbf{7. 多选} MDK 进入 debug 模式后常用操作有：\par
A. 单步调试\par
B. 添加变量到观察窗，查看变量\par
C. 打断点\par
D. 暂停查看程序停在哪\par
\textbf{Ans: A B C D}\par
\textbf{8. 多选} LDO 的工作特点有：\par
A. 效率低\par
B. 输出功率大\par
C. 输出纹波特性好\par
D. 降压幅度大\par
\textbf{Ans: A C}\par
\textbf{9. 单选} \texttt{\#define N 2}，\texttt{\#define Y(n) ((N+1)*n)}，执行 \texttt{z=2*(N+Y(5));} 后：\par
A. z 无定值\par
B. z=70\par
C. 语句有错误\par
D. z=34\par
\textbf{Ans: D}\par
\textbf{10. 判断} 软件定时准确、不需外加硬件电路、没有什么缺点。\par
A. 正确\par
B. 错误\par
\textbf{Ans: B}。软件定时占 CPU，精度受中断/分支/优化影响。\par
\textbf{11. 单选} GPIO 配置成输出时，单个引脚最大输出电流为（）mA，所有引脚总体输出电流不能超过（）mA。\par
A. 10 / 100\par
B. 50 / 500\par
C. 20 / 150\par
D. 10 / 500\par
\textbf{Ans: C}\par
\textbf{12. 单选} C 语言规定，在一个源程序中，\texttt{main} 主函数的位置是：\par
A. 必须在最前面\par
B. 可以在任意位置\par
C. 必须在系统调用的库函数后面\par
D. 必须在最后面\par
\textbf{Ans: B}\par
\textbf{13. 判断} 电路存在共模噪声时，可使用电池供电快速排除电源问题。\par
A. 正确\par
B. 错误\par
\textbf{Ans: A}\par
\textbf{14. 多选} 有关电路原理图，说法正确的有：\par
A. 原理图是电子元器件的支撑体，是电气连接的提供者\par
B. 原理图直接体现电路结构和工作原理，一般用于设计、分析电路\par
C. 识别原理图中的元件符号及连接方式，可了解电路实际工作的潜在情况\par
D. 原理图是用约定符号绘制的一种表示电路结构的图形\par
\textbf{Ans: B C D}。A 描述的是 PCB。\par
\textbf{15. 判断} SYSCLK 时钟源有三个来源：HSI RC、HSE OSC、PLL。\par
A. 正确\par
B. 错误\par
\textbf{Ans: A}\par
\textbf{16. 判断} 端口输入数据寄存器 IDR 的复位值为 \texttt{00000000H}。\par
A. 错误\par
B. 正确\par
\textbf{Ans: A}。IDR 读实际引脚电平，复位后不应机械记作固定 0。\par
\textbf{17. 单选} 下面哪一种不是 STM32 GPIO 的输出模式？\par
A. 浮空输出\par
B. 开漏输出\par
C. 推挽输出\par
D. 复用推挽输出\par
\textbf{Ans: A}\par
\textbf{18. 判断} 向量中断控制器允许有相同的优先级。\par
A. 错误\par
B. 正确\par
\textbf{Ans: B}。允许配置相同优先级；但相同抢占优先级不能互相嵌套打断。\par
\textbf{19. 单选} 三级流水线技术的顺序是：\par
A. 译码 -> 取指 -> 执行\par
B. 取指 -> 译码 -> 执行\par
C. 取指 -> 编译 -> 执行\par
D. 编译 -> 取指 -> 执行\par
\textbf{Ans: B}\par
\textbf{20. 判断} STM32F1 SPI1 的最高频率为 36MHz，即 \texttt{fpclk/2}。\par
A. 正确\par
B. 错误\par
\textbf{Ans: A}。SPI1 挂 APB2；SPI2/3 挂 APB1，最高 18Mbit/s。\par
\textbf{21. 判断} LSE 一般为 32.768kHz，LSI 约 40kHz，可提供看门狗和自动唤醒单元时钟源。\par
A. 错误\par
B. 正确\par
\textbf{Ans: B}\par
\textbf{22. 判断} Cortex-M3 内核采用冯·诺依曼结构的三级流水线。\par
A. 错误\par
B. 正确\par
\textbf{Ans: A}。Cortex-M3 是哈佛结构 + 三级流水线。\par
\textbf{23. 单选} ADC 转换过程正确的是：\par
A. 编码 -> 采样 -> 量化\par
B. 量化 -> 采样 -> 编码\par
C. 采样 -> 量化 -> 编码\par
D. 采样 -> 编码 -> 量化\par
\textbf{Ans: C}\par
\textbf{24. 判断} PWM1 和 PWM2 模式的区别仅在于相位 180 度，因此没有意义。\par
A. 正确\par
B. 错误\par
\textbf{Ans: B}\par
\textbf{25. 多选} 关于 STM32 PWM，说法正确的有：\par
A. PWM 输出占空比由 \texttt{TIMx\_CCR} 决定\par
B. PWM 又称脉冲宽度调制\par
C. PWM 输出频率由 \texttt{TIMx\_ARR} 决定\par
D. STM32 的通用和高级定时器都可以产生 PWM\par
\textbf{Ans: A B C D}\par
\textbf{26. 判断} CMSIS 是强制使用的芯片设计标准，且 \texttt{GPIO\_ResetBits()} 属于 CMSIS。\par
A. 错误\par
B. 正确\par
\textbf{Ans: A}。CMSIS 是接口标准；\texttt{GPIO\_ResetBits()} 属于 ST 标准外设库。\par
\textbf{27. 判断} \texttt{stm32f10x.h} 相当于 \texttt{reg52.h}；\texttt{stm32f10x\_conf.h} 负责控制外设库头文件开关。\par
A. 错误\par
B. 正确\par
\textbf{Ans: B}\par
\textbf{28. 判断} STM32 NVIC 中固定优先级最高的是 -3。\par
A. 正确\par
B. 错误\par
\textbf{Ans: A}。Reset=-3，NMI=-2，HardFault=-1。\par
\textbf{29. 判断} STM32F1 复位后所有 I/O 口均为浮空输入状态。\par
A. 错误\par
B. 正确\par
\textbf{Ans: B}。考试题通常按普通 GPIO 记忆；JTAG/SWD 引脚有默认复用例外。\par
\textbf{30. 单选} 在 STM32 中，如何通过 C 语言访问 GPIOB 的 ODR 以设置引脚输出高电平？\par
A. 通过串口发送指令控制 GPIOB\par
B. 调用 \texttt{GPIO\_SetBits(GPIOB, 0xFFFF);}\par
C. 使用指针访问：\texttt{*(unsigned int*)(0x40010C0C) = 0xFFFF;}\par
D. 直接调用 \texttt{GPIOB-\textgreater{}ODR = 0xFFFF;} 无需定义\par
\textbf{Ans: C}\par
\subsection*{17.2 新增覆盖题}
\textbf{31. 单选} STM32F103 复位后默认系统时钟源通常是：\par
A. HSI 8MHz\par
B. HSE 8MHz\par
C. PLL 72MHz\par
D. LSE 32.768kHz\par
\textbf{Ans: A}\par
\textbf{32. 多选} STM32F1 中挂在 APB2 总线上的外设通常包括：\par
A. GPIOA-G\par
B. USART1\par
C. SPI1\par
D. TIM2\par
\textbf{Ans: A B C}。TIM2 在 APB1。\par
\textbf{33. 单选} \texttt{BOOT0=1, BOOT1=0} 时，STM32F1 从哪里启动？\par
A. 主 Flash\par
B. 系统存储器 Bootloader\par
C. SRAM\par
D. 外部 EEPROM\par
\textbf{Ans: B}\par
\textbf{34. 多选} 配置 EXTI0 外部中断通常需要：\par
A. 使能 GPIO 和 AFIO 时钟\par
B. 用 \texttt{GPIO\_EXTILineConfig()} 选择 PA0/PB0/... 映射到 EXTI0\par
C. 配置 \texttt{EXTI\_Init()} 的触发边沿\par
D. 配置 NVIC 并编写 \texttt{EXTI0\_IRQHandler()}\par
\textbf{Ans: A B C D}\par
\textbf{35. 单选} EXTI pending 标志通常如何清除？\par
A. 读 \texttt{EXTI\_PR} 自动清除\par
B. 对 \texttt{EXTI\_PR} 对应位写 0 清除\par
C. 对 \texttt{EXTI\_PR} 对应位写 1 清除\par
D. 关闭 RCC 时钟清除\par
\textbf{Ans: C}\par
\textbf{36. 单选} \texttt{GPIOx\_BSRR} 的高 16 位写 1 的作用是：\par
A. 设置对应 ODR 位为 1\par
B. 复位对应 ODR 位为 0\par
C. 读取对应 IDR 位\par
D. 锁定 GPIO 配置\par
\textbf{Ans: B}\par
\textbf{37. 多选} 关于 GPIO 开漏输出，正确的是：\par
A. 只能主动拉低，输出高电平依赖上拉\par
B. 适合 I2C 等线与总线\par
C. 与推挽输出一样内部上下管都主动驱动\par
D. 复用开漏输出也常需要上拉电阻\par
\textbf{Ans: A B D}\par
\textbf{38. 单选} SysTick 是几位定时器？计数方向如何？\par
A. 16 位向上计数\par
B. 16 位向下计数\par
C. 24 位向下计数\par
D. 32 位向上计数\par
\textbf{Ans: C}\par
\textbf{39. 单选} \texttt{SysTick\_Config(SystemCoreClock/1000)} 在 \texttt{SystemCoreClock=72MHz} 时通常产生：\par
A. 1us 中断\par
B. 1ms 中断\par
C. 10ms 中断\par
D. 1s 中断\par
\textbf{Ans: B}\par
\textbf{40. 多选} USART 发送/接收相关状态位正确的是：\par
A. \texttt{TXE} 表示发送数据寄存器空\par
B. \texttt{TC} 表示一帧发送完成\par
C. \texttt{RXNE} 表示接收数据寄存器非空\par
D. \texttt{ORE} 表示溢出错误\par
\textbf{Ans: A B C D}\par
\textbf{41. 单选} USART1 的 TX/RX 引脚在常用默认映射下分别是：\par
A. PA9 / PA10\par
B. PA2 / PA3\par
C. PB6 / PB7\par
D. PA5 / PA6\par
\textbf{Ans: A}\par
\textbf{42. 多选} SPI 与 I2C 的典型区别正确的是：\par
A. SPI 通常全双工，I2C 通常半双工\par
B. SPI 有 SCK/MOSI/MISO/NSS，I2C 有 SCL/SDA\par
C. I2C 需要开漏和上拉电阻\par
D. SPI 通过地址选择从设备，I2C 只靠片选线选择从设备\par
\textbf{Ans: A B C}\par
\textbf{43. 单选} I2C 总线的起始条件是：\par
A. SCL 高电平时 SDA 由高变低\par
B. SCL 高电平时 SDA 由低变高\par
C. SCL 低电平时 SDA 由高变低\par
D. SCL 低电平时 SDA 由低变高\par
\textbf{Ans: A}\par
\textbf{44. 单选} 定时器定时时间公式正确的是：\par
A. \texttt{T=(PSC+1)(ARR+1)/TIMxCLK}\par
B. \texttt{T=TIMxCLK/((PSC+1)(ARR+1))}\par
C. \texttt{T=PSC*ARR*TIMxCLK}\par
D. \texttt{T=ARR/(PSC+TIMxCLK)}\par
\textbf{Ans: A}\par
\textbf{45. 单选} \texttt{TIMx\_PSC=71, TIMx\_ARR=999, TIMxCLK=72MHz} 时更新频率约为：\par
A. 1Hz\par
B. 10Hz\par
C. 1kHz\par
D. 72kHz\par
\textbf{Ans: C}。周期 \texttt{72*1000/72MHz=1ms}。\par
\textbf{46. 多选} 关于 STM32F1 定时器，正确的是：\par
A. TIM6/TIM7 是基本定时器，无 GPIO 输出通道\par
B. TIM2-TIM5 是通用定时器\par
C. TIM1/TIM8 是高级定时器，可支持互补输出、死区、刹车\par
D. 所有定时器都能直接产生 PWM 输出\par
\textbf{Ans: A B C}\par
\textbf{47. 单选} ADC 输入引脚应配置成什么 GPIO 模式？\par
A. 推挽输出\par
B. 浮空输入\par
C. 模拟输入\par
D. 复用推挽输出\par
\textbf{Ans: C}\par
\textbf{48. 多选} STM32 ADC 相关说法正确的是：\par
A. F1 ADC 通常为 12 位 SAR ADC\par
B. 转换时间 = 采样时间 + 12.5 个 ADC 时钟周期\par
C. ADCCLK 不能超过规定上限，F1 常记最大 14MHz\par
D. ADC 可配合 DMA 将转换结果搬运到内存\par
\textbf{Ans: A B C D}\par
\textbf{49. 单选} DMA 循环模式 Circular 的特点是：\par
A. 传输一次后自动关闭通道且不能再启用\par
B. 到达缓冲区末尾后自动回到开头继续传输\par
C. 只能用于内存到内存传输\par
D. 只能用于 DAC，不能用于 ADC/USART\par
\textbf{Ans: B}\par
\textbf{50. 多选} DMA 配置中常见字段及含义正确的是：\par
A. \texttt{DMA\_CPARx} 为外设地址\par
B. \texttt{DMA\_CMARx} 为存储器地址\par
C. \texttt{DMA\_CNDTRx} 为传输数量\par
D. \texttt{DMA\_CCRx} 中可设置方向、循环、地址递增、数据宽度、优先级\par
\textbf{Ans: A B C D}\par
\textbf{51. 单选} Cortex-M3 向量表第 0 项和第 1 项分别存放：\par
A. Reset\_Handler 和 HardFault\_Handler\par
B. 初始 MSP 和 Reset\_Handler\par
C. 初始 PC 和初始 LR\par
D. SysTick\_Handler 和 PendSV\_Handler\par
\textbf{Ans: B}\par
\textbf{52. 多选} 下列哪些通常属于“写 1 清除”类标志？\par
A. \texttt{EXTI\_PR} 对应 pending 位\par
B. \texttt{DMA\_IFCR} 中对应清除位\par
C. \texttt{GPIOx\_IDR} 输入数据位\par
D. \texttt{GPIOx\_BSRR} 的置位/复位控制位\par
\textbf{Ans: A B}。D 是写 1 执行动作，不是清状态标志。\par
\textbf{53. 单选} 在 STM32F1 标准外设库中，修改 TIM3\_CH2 的 PWM 占空比通常使用：\par
A. \texttt{TIM\_SetCompare2(TIM3, value)}\par
B. \texttt{GPIO\_SetBits(TIM3, value)}\par
C. \texttt{RCC\_APB2PeriphClockCmd(TIM3, value)}\par
D. \texttt{USART\_SendData(TIM3, value)}\par
\textbf{Ans: A}\par
\textbf{54. 多选} C 语言中 \texttt{volatile} 常用于修饰：\par
A. 中断服务程序与主程序共享的变量\par
B. 硬件寄存器映射地址对应的变量\par
C. 会被外部硬件改变的变量\par
D. 普通局部循环变量，且永远都必须加\par
\textbf{Ans: A B C}\par
\textbf{55. 单选} 若宏定义为 \texttt{\#define SQR(x) x*x}，则 \texttt{SQR(1+2)} 的结果是：\par
A. 9\par
B. 5\par
C. 6\par
D. 编译错误\par
\textbf{Ans: B}。展开为 \texttt{1+2*1+2=5}，正确写法应为 \texttt{\#define SQR(x) ((x)*(x))}。\par
\vspace{1pt}\hrule\vspace{1pt}
\section*{18. 公式与代码片段速备}
\subsection*{18.1 STM32 GPIO 标准库初始化模板}
\begin{Code}
RCC_APB2PeriphClockCmd(RCC_APB2Periph_GPIOA, ENABLE);
GPIO_InitTypeDef gi;
gi.GPIO_Pin   = GPIO_Pin_5;
gi.GPIO_Mode  = GPIO_Mode_Out_PP;        // 推挽输出
gi.GPIO_Speed = GPIO_Speed_50MHz;
GPIO_Init(GPIOA, &gi);
GPIO_SetBits(GPIOA, GPIO_Pin_5);          // 置高
\end{Code}
\subsection*{18.2 SysTick 1 ms}
\begin{Code}
if (SysTick_Config(SystemCoreClock / 1000)) while (1);   // 失败死循环
void SysTick_Handler(void) { tick_ms++; }
\end{Code}
\subsection*{18.3 EXTI 外部中断}
\begin{Code}
RCC_APB2PeriphClockCmd(RCC_APB2Periph_AFIO|RCC_APB2Periph_GPIOA, ENABLE);
GPIO_EXTILineConfig(GPIO_PortSourceGPIOA, GPIO_PinSource0);
EXTI_InitTypeDef ei = {EXTI_Line0, EXTI_Mode_Interrupt, EXTI_Trigger_Rising, ENABLE};
EXTI_Init(&ei);
NVIC_InitTypeDef ni = {EXTI0_IRQn, 1, 1, ENABLE};
NVIC_Init(&ni);
\end{Code}
\subsection*{18.4 PWM (TIM3 CH1)}
\begin{Code}
TIM_TimeBaseInitTypeDef tb = { .TIM_Period=999, .TIM_Prescaler=71,        // 1 kHz @ 72MHz
                               .TIM_CounterMode=TIM_CounterMode_Up };
TIM_TimeBaseInit(TIM3, &tb);
TIM_OCInitTypeDef oc = { .TIM_OCMode=TIM_OCMode_PWM1, .TIM_OutputState=TIM_OutputState_Enable,
                         .TIM_Pulse=500, .TIM_OCPolarity=TIM_OCPolarity_High };
TIM_OC1Init(TIM3, &oc);                                                   // 50% 占空比
TIM_Cmd(TIM3, ENABLE);
\end{Code}
\subsection*{18.5 ADC 单次 + DMA}
\begin{Code}
RCC_ADCCLKConfig(RCC_PCLK2_Div6);                                          // 12 MHz
ADC_InitTypeDef ai = {ADC_Mode_Independent, DISABLE, DISABLE,
                      ADC_ExternalTrigConv_None, ADC_DataAlign_Right, 1};
ADC_Init(ADC1, &ai);
ADC_RegularChannelConfig(ADC1, ADC_Channel_0, 1, ADC_SampleTime_55Cycles5);
ADC_DMACmd(ADC1, ENABLE);
ADC_Cmd(ADC1, ENABLE);
ADC_ResetCalibration(ADC1); while(ADC_GetResetCalibrationStatus(ADC1));
ADC_StartCalibration(ADC1);  while(ADC_GetCalibrationStatus(ADC1));
ADC_SoftwareStartConvCmd(ADC1, ENABLE);
\end{Code}
\vspace{1pt}\hrule\vspace{1pt}
\hint{临场提示：① 看到题目说"哈佛 / 冯·诺依曼"先想 Cortex‑M3 = 哈佛、8051 = 哈佛、Cortex‑A 系列 = 冯诺依曼简化扩展。}\par
\hint{② 见到 PWM 立刻条件反射：\textbf{ARR↔频率，CCR↔占空比}。}\par
\hint{③ 见到 SPI 频率立刻反应：\textbf{fPCLK/2}；F103 中 SPI1 上限 \textbf{36 Mbit/s}，SPI2/3 上限 \textbf{18 Mbit/s}。}\par
\hint{④ 见到 ADC 流程：\textbf{采—量—编}。}\par
\hint{⑤ 看到 \texttt{\#define} 立刻数括号：宏展开 = 文本替换。}\par
\end{multicols*}
\end{document}
